\subsection{The masked fold chain and confidential local checks}\label{sec:full:fold-oracle}

\begin{definition}[Honest masked fold-oracle chain]\label{def:full:masked-oracles}
Before the evaluation point $z$ is known, sample
$R\overset{\$}{\leftarrow}\F[X]_{<B}$ and set
\[
 Q_0=U_f^{(m)}+X^NR.
\]
After the sumcheck round challenge $r_j$,
define $Q_{j+1}=\pfold_{r_j}(Q_j)$
for $0\le j<m$. At each level set
\[
 w_j(x):=Q_j(x)\quad(x\in\Omega_j),
\]
and commit to \emph{every} $w_j(x)$ with
fresh independent field-symbol commitment
$\mathsf{VC}(w_j(x);\mathbf r_{j,x})$.
The resulting Merkle root $\mathsf{root}_j$
is published \emph{before} that level's
local authentication-query indices
are selected.
\end{definition}

\begin{lemma}[Exact degree and word-fold invariants]\label{lem:full:fold-invariants}
For an honest masked chain,
\[
 \deg Q_j<D_j,\qquad
 w_{j+1}=\mathsf{WFold}_{r_j,j}(w_j)
\]
for every $0\le j<m$.
At the final level,
$Q_m=a+\rho X$ with
$a=\widetilde f(r)$.
\end{lemma}
\begin{proof}
The degree and codeword statements
are respectively
Theorem~\ref{thm:kbff:code-rate}
and Corollary~\ref{cor:kbff:wordchain}.
The final linear form and its
constant coefficient follow from
Theorem~\ref{thm:kbff:linear-closure}.
\end{proof}

\begin{definition}[Confidential authenticated fold query]\label{def:full:fold-query}
For a query $(j,y)$ with
$0\le j<m$ and $y\in\Omega_{j+1}$,
choose either public square root
$x\in\Omega_j$ satisfying $x^2=y$.
The prover gives public Merkle
authentication shells for
$(j,x)$, $(j,-x)$ and $(j+1,y)$.
It then proves in zero knowledge
that the three authenticated
commitment vectors possess
\emph{valid field-symbol openings}
$(a_+,a_-,b)$ satisfying
\[
 \boxed{
 b=(2r_j-1)\frac{a_++a_-}{2}
   +(1-r_j)\frac{a_+-a_-}{2x}.}
\]
The proof uses the same hidden
openings as the corresponding
registered word, not independently
chosen values for each query.
No $a_+,a_-,b$ or commitment
blindings are disclosed.
\end{definition}

\begin{lemma}[The query equation is independent of square-root choice]\label{lem:full:fold-query-choice}
Replacing $x$ by $-x$ and
interchanging the two level-$j$
openings leaves the equation
of Definition~\ref{def:full:fold-query}
unchanged.
\end{lemma}
\begin{proof}
The sum $a_++a_-$ is unchanged.
The difference $a_+-a_-$ and
the denominator $2x$ both
change sign, so their quotient
is unchanged. The coefficients
$(2r_j-1)$ and $(1-r_j)$
do not depend on the root.
\end{proof}

\begin{lemma}[Honest confidential fold queries pass]\label{lem:full:fold-query-complete}
Every honest authenticated fold query
has valid committed field openings
satisfying its hidden fold equation.
\end{lemma}
\begin{proof}
By Definition~\ref{def:full:masked-oracles},
the three openings are the
true values
$Q_j(x),Q_j(-x),Q_{j+1}(y)$.
The word-fold identity of
Lemma~\ref{lem:full:fold-invariants}
is exactly the query equation
at $y=x^2$. The blinding
values are those generated
for the registered leaves.
Therefore all queried statements
have valid witnesses.
\end{proof}

\begin{definition}[Confidential final closure]\label{def:full:final-close}
Let $Q_m(X)=a+\rho X$ be the
last hidden polynomial.
The public verifier requests
authenticated commitment shells
at two publicly fixed, distinct
points $x_0,x_1\in\Omega_m$.
The prover supplies a zero-knowledge
relation proof that:
\begin{enumerate}[label=\textup{(\roman*)}]
\item the two hidden authenticated
field values $u_0,u_1$ lie on
a common degree-$<2$ polynomial;
\item the constant $a$ of that
line satisfies
$x_1u_0-x_0u_1=(x_1-x_0)a$;
\item $a$ agrees with the hidden
sumcheck terminal relation
$\mu_m(r_{m-1})=\chi_m a+b_m$.
\end{enumerate}
The proof binds $a$ to the
same final degree-$<2$ word
as the registered final root.
\end{definition}

\begin{lemma}[Honest final closure is valid]\label{lem:full:final-complete}
The final closure relation of
Definition~\ref{def:full:final-close}
has a witness in every honest
execution, provided $|\Omega_m|\ge2$.
\end{lemma}
\begin{proof}
The group domain has
$M_m=2^s\ge2$ elements by
Lemma~\ref{lem:kbff:domain-sizes},
so distinct $x_0,x_1$ exist.
By Lemma~\ref{lem:full:fold-invariants},
$Q_m=a+\rho X$. Its two
evaluations lie on this line
and Lemma~\ref{lem:kbff:two-point-line}
gives the required constant
coefficient equation.
Corollary~\ref{cor:full:masked-final-link}
gives the final hidden sumcheck
equation, proving every item.
\end{proof}

\begin{lemma}[Degree-certified final word is uniquely determined by two values]\label{lem:full:two-unique}
If $P,Q\in\F[X]_{<2}$ agree at
two distinct points of $\F$,
then $P=Q$.
\end{lemma}
\begin{proof}
The difference $P-Q$ has
degree at most one and
two distinct roots.
The univariate root bound
Lemma~\ref{lem:history:rootcount}
forces $P-Q=0$.
\end{proof}

\begin{remark}[The proximity gap]
A zero-knowledge proof of the
\emph{entire} degree-$<D_j$ word
and each fold relation is a valid
but nonsuccinct baseline.
Replacing that full relation
by sparse local checks requires
a proximity theorem for
\emph{arbitrary adversarial words}
in the enlarged degree budgets
$D_j$, not merely a root-count
lemma for two known polynomials.
Chapter~14 keeps those two
verification modes separate.
\end{remark}
